\documentclass[11pt]{article}
\usepackage[margin=1in]{geometry}
\usepackage{amsmath,amssymb,amsfonts,mathtools}
\usepackage{array,booktabs,enumitem,graphicx,xcolor,hyperref}
\usepackage[T1]{fontenc}
\usepackage[utf8]{inputenc}
\title{On the connection between Quantum Mechanics and the geometry of
  two-dimensional strings}
\author{Source-backed arXiv sample}
\date{}
\begin{document}
\maketitle
\begin{abstract}
On the basis of an area-preserving symmetry in the phase space of a one-dimensional matrix model - believed to describe two-dimensional string theory in a black-hole background which also allows for space-time foam - we give a geometric interpretation of the fact that two-dimensional stringy black holes are consistent with conventional quantum mechanics due to the infinite gauged `W-hair' property that characterises them.
\end{abstract}
\section{References And Citations}
\label{sec:refs}
This sample cites source keys \cite{Ell} and \cite{Gro}. Section~\ref{sec:refs} and Equation~\ref{eq:source-demo} provide cross-reference coverage.
\begin{equation}
\label{eq:source-demo}
a^2 + b^2 = c^2
\end{equation}
\begin{thebibliography}{9}
\bibitem{Ell} Source bibliography entry 1 extracted from arXiv metadata/source key `Ell`.
\bibitem{Gro} Source bibliography entry 2 extracted from arXiv metadata/source key `Gro`.
\bibitem{Ant} Source bibliography entry 3 extracted from arXiv metadata/source key `Ant`.
\bibitem{Witt} Source bibliography entry 4 extracted from arXiv metadata/source key `Witt`.
\end{thebibliography}

\end{document}
